\chapter{In Purificatione sanctae Mariae}

\rubrica{Ad completorium psalmi consueti:\quickref{psalmi}
  quos sequitur antiphona}\edissue{sequitur antiphona -- this may or may not mean arrangement different from the usual super psalmos antiphona. Look around for analogical rubrics elsewhere. Is it related to the feast's -- presumably lower -- rank? Or with the pre-Lenten period?}
Verbum caro factum.\quickref{form:nativitas}
\rubrica{Capitulum}
Parvulus natus.\quickref{form:nativitas}
\rubrica{Hymnus.}
Corde natus.\quickref{hym:corde}
\rubrica{V.} Benedictus qui venit.\quickref{nativitas:verse:benedictus}
\rubrica{Ad Nunc dimittis antiphona.}
Magnum nomen Domini.\quickref{nativitas:magnum}
\rubrica{Haec omnia require superius in completorio
  de Nativitate Domini.\quickref{form:nativitas}}

\rubrica{Preces solitae:\quickref{preces}
  adiuncto in fine versiculo isto.}\\
Post partum \editorial{Virgo inviolata permansisti.
  Dei Genitrix intercede pro nobis}.\footcite[De BMV, 3v]{diu1523}
\rubrica{Oratio.}
Deus qui sa\editorial{lutis.}\quickref{circumcisio:oratio}
\rubrica{ut supra in Circumcisione Domini:
  et haec sola.}
